\chapter{Konstrukcje cyrklem i linijką}

\section{Które operacje są dozwolone?}

Zaczynamy od punktów $O=(0,0)$ i $I=(1,0)$, ustalających jednostkę. Wolno prowadzić prostą przez dwa różne znane punkty, rysować okrąg o znanym środku i promieniu równym odległości dwóch znanych punktów, oraz wybierać izolowane punkty przecięcia tak otrzymanych prostych i okręgów. Linijka nie ma podziałki. Pokrywające się proste lub okręgi nie pozwalają wybierać dowolnego nowego punktu.

Punkt jest \emph{konstruowalny}, jeśli powstaje po skończenie wielu krokach. Liczbę rzeczywistą $a$ nazywamy konstruowalną, jeśli konstruowalny jest punkt $(a,0)$. Przez prostopadłe i przenoszenie odcinków równoważnie możemy mówić o konstruowalności obu współrzędnych punktu $(a,b)$.

Przypomnijmy podstawowe operacje geometryczne. Przecięcia dwóch okręgów o równych promieniach wyznaczają symetralną odcinka łączącego ich środki: równość kwadratów odległości po odjęciu daje właśnie równanie tej prostej. Pozwala to konstruować środki odcinków i prostopadłe, a przez dwie kolejne prostopadłe także proste równoległe. Dodawanie odcinków, odbijanie ich względem początku i przenoszenie cyrklem daje dodawanie, odejmowanie oraz zmianę znaku współrzędnej.

Dokładniej, prostopadłą przez punkt $P$ do danej prostej otrzymujemy jako symetralną cięciwy wyciętej na niej przez okrąg o środku $P$. Promień wybieramy większy niż odległość $P$ od prostej; można go uzyskać przez kolejne podwajanie jednostki. Końce cięciwy są równoodległe od $P$, więc jej symetralna przechodzi przez $P$.

\section{Dokładne kryterium algebraiczne}

\begin{twierdzenie}[Wieże kwadratowe]
Liczba rzeczywista $a$ jest konstruowalna wtedy i tylko wtedy, gdy należy do pewnego ciała $K_r$ w wieży rzeczywistych rozszerzeń
\[
\Q=K_0\subseteq K_1\subseteq\cdots\subseteq K_r,\qquad
K_i=K_{i-1}(\sqrt{b_i}),\quad b_i\in K_{i-1},\ b_i>0.
\]
Kroki, które rzeczywiście powiększają ciało, mają stopień $2$.
\end{twierdzenie}
\begin{proof}
\emph{Od konstrukcji do wieży.} Załóżmy, że dotychczasowe współrzędne należą do rzeczywistego ciała $K$. Prosta przez dwa znane punkty ma równanie liniowe ze współczynnikami w $K$. Okrąg ma równanie
\[
(x-u)^2+(y-v)^2=p,\qquad u,v,p\in K,
\]
bo kwadrat promienia jest sumą kwadratów różnic znanych współrzędnych. Przecięcie dwóch nierównoległych prostych otrzymujemy przez rozwiązanie układu liniowego, więc jego współrzędne są w $K$.

Przy przecięciu prostej z okręgiem parametryzujemy prostą w postaci $(x,y)=(x_0,y_0)+t(v_1,v_2)$ z parametrami w $K$. Po podstawieniu otrzymujemy równanie kwadratowe w $t$. Jego rzeczywiste rozwiązania należą do $K$ albo do $K(\sqrt b)$ dla dodatniego $b\in K$. Przy dwóch okręgach o różnych środkach odejmowanie ich równań daje prostą, więc sytuacja redukuje się do poprzedniej. Okręgi współśrodkowe są rozłączne albo pokrywają się i nie dają izolowanych nowych punktów. Każdy dozwolony krok powiększa zatem ciało co najwyżej kwadratowo. Indukcja po liczbie kroków daje wieżę z tezy.

\emph{Od wieży do konstrukcji.} Trzeba wykazać, że konstruowalne liczby tworzą ciało i można wyciągać z dodatnich takich liczb pierwiastki kwadratowe. Dodawanie i odejmowanie opisaliśmy wyżej. Dla mnożenia weźmy prostą przez $(1,0)$ i $(0,a)$. Ma równanie $y=a(1-x)$. Równoległa do niej przez $(b,0)$ przecina oś pionową w $(0,ab)$, co konstruuje iloczyn. Dla $a\ne0$ równoległa do tej samej prostej przez $(0,b)$ przecina oś poziomą w $(b/a,0)$. Przypadek mnożenia przez zero jest natychmiastowy. Otrzymujemy więc wszystkie działania ciała, w tym liczby wymierne.

Dla konstruowalnego $b>0$ rysujemy okrąg o średnicy od $(-1,0)$ do $(b,0)$. Środek i promień są konstruowalne. Przecięcie z osią pionową spełnia
\[
\left(\frac{b-1}{2}\right)^2+y^2=\left(\frac{b+1}{2}\right)^2,
\qquad y^2=b.
\]
Górne przecięcie daje $\sqrt b$. Indukcyjnie wszystkie elementy każdego $K_i$ są więc konstruowalne, a w szczególności $a$.

Każde rzeczywiste rozszerzenie stopnia $2$ ma używaną wyżej postać: jeżeli generuje je $\alpha$ o wielomianie minimalnym $T^2+pT+q$, to $(2\alpha+p)^2=p^2-4q$. Wyróżnik jest dodatni, bo $\alpha$ jest rzeczywiste, a nierozkładalny wielomian w charakterystyce zero nie ma zera podwójnego. Zatem ciało jest generowane przez pierwiastek kwadratowy z wyróżnika.
\end{proof}

\begin{wniosek}
Jeżeli $a$ jest konstruowalne, to jest algebraiczne nad $\Q$, a stopień $[\Q(a):\Q]$ jest potęgą liczby $2$.
\end{wniosek}
\begin{proof}
Po usunięciu trywialnych kroków wieża z twierdzenia ma stopień $2^r$ nad $\Q$. Ponieważ $\Q\subseteq\Q(a)\subseteq K_r$, twierdzenie o wieży daje $[\Q(a):\Q]\mid 2^r$. Każdy dodatni dzielnik $2^r$ jest potęgą $2$.
\end{proof}

To warunek konieczny. Pełnym kryterium jest przynależność do wieży kwadratowej; samo stwierdzenie, że wielomian minimalny ma stopień $4$, $8$ lub $16$, nie zastępuje zbudowania takiej wieży.

\section{Podwojenie sześcianu i trysekcja kąta}

\begin{twierdzenie}
Nie można cyrklem i linijką podwoić objętości danego sześcianu ani dokonać trysekcji dowolnego danego kąta.
\end{twierdzenie}
\begin{proof}
\emph{Sześcian.} Przy boku początkowym $1$ nowy bok musiałby wynosić $\sqrt[3]{2}$. Jego wielomian minimalny to $T^3-2$, którego nierozkładalność wykazano w przykładzie 2.7. Stopień $3$ nie jest potęgą $2$, więc stosujemy wniosek 6.2.

\emph{Kąt.} Kąt $60^\circ$ jest konstruowalny, na przykład jako kąt trójkąta równobocznego. Gdyby zawsze można było dokonać trysekcji, można byłoby skonstruować kąt $20^\circ$, a przez przecięcie jego ramienia z okręgiem jednostkowym również $u=2\cos20^\circ$. Tożsamość $2\cos3\theta=(2\cos\theta)^3-3(2\cos\theta)$ daje
\[
u^3-3u-1=0.
\]
Wymierne zero unormowanego wielomianu całkowitego musi być całkowite i dzielić wyraz wolny: wynika to po podstawieniu nieskracalnego ułamka i porównaniu podzielności. Możliwe byłyby tu tylko $1$ i $-1$, ale wartości wielomianu to odpowiednio $-3$ i $1$. Nie ma więc zera wymiernego. Wielomian stopnia $3$ jest nierozkładalny, a $u$ ma stopień $3$, ponownie niemożliwy dla liczby konstruowalnej.
\end{proof}

Niemożliwa jest \emph{ogólna} trysekcja. Niektóre kąty można podzielić na trzy równe części, na przykład $90^\circ$ na trzy kąty po $30^\circ$. Twierdzenie nie dotyczy też konstrukcji z linijką z podziałką ani origami: są to inne zestawy dozwolonych operacji.

\section{Wielokąty foremne: dwa pełne przykłady}

Foremny $n$-kąt wpisany w okrąg jednostkowy można skonstruować, jeśli można skonstruować punkt $(\cos(2\pi/n),\sin(2\pi/n))$. Kolejne wierzchołki otrzymujemy przez odkładanie tej samej cięciwy. Odwrotnie, z gotowego wielokąta i jego środka otrzymujemy taki punkt po przeniesieniu i przeskalowaniu figury.

\begin{przyklad}[Pięciokąt]
Niech $\zeta=e^{2\pi i/5}$ i $u=\zeta+\zeta^{-1}=2\cos(2\pi/5)>0$. Suma szeregu geometrycznego daje $1+\zeta+\cdots+\zeta^4=0$. Dzieląc przez $\zeta^2$ i używając $\zeta^2+\zeta^{-2}=u^2-2$, otrzymujemy $u^2+u-1=0$. Stąd
\[
\cos\frac{2\pi}{5}=\frac{\sqrt5-1}{4},\qquad
\sin\frac{2\pi}{5}=\sqrt{1-\left(\frac{\sqrt5-1}{4}\right)^2}.
\]
Obie współrzędne uzyskujemy przez działania ciała i dodatnie pierwiastki kwadratowe. Twierdzenie 6.1 daje konstrukcję pięciokąta.
\end{przyklad}

\begin{przyklad}[Siedmiokąt]
Teraz niech $\zeta=e^{2\pi i/7}$ i $u=\zeta+\zeta^{-1}=2\cos(2\pi/7)$. Po podzieleniu $1+\zeta+\cdots+\zeta^6=0$ przez $\zeta^3$ dostajemy
\[
1+u+(u^2-2)+(u^3-3u)=0,\qquad u^3+u^2-2u-1=0.
\]
Jedynymi kandydatami na zera wymierne są $\pm1$, lecz wartości wynoszą odpowiednio $-1$ i $1$. Wielomian jest nierozkładalny stopnia $3$. Gdyby siedmiokąt był konstruowalny, konstruowalne byłoby $u$, co przeczy wnioskowi 6.2. Siedmiokąta foremnego nie można więc skonstruować cyrklem i linijką.
\end{przyklad}

\subsection*{Wielomiany cyklotomiczne i pełna klasyfikacja}

Niech $\zeta_n=e^{2\pi i/n}$ i niech $\varphi(n)$ oznacza liczbę klas odwracalnych modulo $n$. Wielomian cyklotomiczny definiujemy przez
\[
\Phi_n(T)=\prod_{\substack{1\le a\le n\\\gcd(a,n)=1}}(T-\zeta_n^a).
\]
Współczynniki $\Phi_n$ są całkowite. Istotnie, pierwiastki wielomianu $T^n-1$ dzielą się na rozłączne zbiory pierwiastków pierwotnych rzędów $d\mid n$, więc
$T^n-1=\prod_{d\mid n}\Phi_d(T)$. Indukcja po $n$ i dzielenie wielomianów unormowanych w $\Z[T]$ pokazują, że $\Phi_n\in\Z[T]$.

\begin{lemat}[Nierozkładalność wielomianów cyklotomicznych]
Dla każdego $n\ge1$ wielomian $\Phi_n$ jest nierozkładalny nad $\Q$. W konsekwencji
$[\Q(\zeta_n):\Q]=\varphi(n)$ oraz
$\Gal(\Q(\zeta_n)/\Q)\simeq(\Z/n\Z)^\times$.
\end{lemat}
\begin{proof}
Niech $f\in\Z[T]$ będzie unormowanym wielomianem minimalnym $\zeta_n$ nad $\Q$. Dzieli on $\Phi_n$, na mocy lematu Gaussa. Pokażemy, że jeśli $\ell$ jest liczbą pierwszą i $\ell\nmid n$, to $\zeta_n^\ell$ także zeruje $f$.

Przypuśćmy przeciwnie i oznaczmy przez $g\in\Z[T]$ unormowany wielomian minimalny $\zeta_n^\ell$. Jest to inny nierozkładalny czynnik $\Phi_n$. Ponieważ $g(\zeta_n^\ell)=0$, mamy $f(T)\mid g(T^\ell)$ w $\Z[T]$. Po redukcji modulo $\ell$ zachodzi
$\overline{g(T^\ell)}=\overline{g(T)}^\ell$; zatem $\bar f$ i $\bar g$ mają wspólny nierozkładalny czynnik. Z drugiej strony $fg\mid T^n-1$, także po redukcji, a $T^n-1$ nad $\F_\ell$ jest bez pierwiastków wielokrotnych, gdyż jego pochodna $nT^{n-1}$ jest z nim względnie pierwsza. Sprzeczność.

Ten sam argument stosuje się do dowolnego pierwotnego pierwiastka $n$-tego stopnia. Rozkładając każde $a$ względnie pierwsze z $n$ na czynniki pierwsze, widzimy zatem, że $f(\zeta_n^a)=0$ dla wszystkich takich $a$. Stąd $\Phi_n\mid f$, więc $f=\Phi_n$. Liczba różnych pierwiastków wynosi $\varphi(n)$, co daje stopień. Każdy automorfizm wysyła $\zeta_n$ na $\zeta_n^a$ z $\gcd(a,n)=1$, a wszystkie takie obrazy występują jako sprzężenia. Otrzymujemy wskazany izomorfizm grup.
\end{proof}

\begin{twierdzenie}[Gaussa--Wantzela]
Dla $n\ge3$ foremny $n$-kąt jest konstruowalny cyrklem i linijką wtedy i tylko wtedy, gdy
\[
n=2^r p_1\cdots p_s,\qquad r\ge0,
\]
gdzie $p_1,\ldots,p_s$ są parami różnymi liczbami pierwszymi Fermata, czyli liczbami pierwszymi postaci $2^{2^m}+1$.
\end{twierdzenie}
\begin{proof}
Połóżmy $u=\zeta_n+\zeta_n^{-1}=2\cos(2\pi/n)$. Konstrukcja wielokąta jest równoważna konstruowalności $u$: w jedną stronę odczytujemy współrzędną wierzchołka, a w drugą konstruujemy $\cos(2\pi/n)=u/2$ i dodatnią liczbę $\sin(2\pi/n)=\sqrt{1-u^2/4}$ (dla $n\ge3$). Następnie odkładamy odpowiednią cięciwę.

Jeżeli wielokąt jest konstruowalny, to $u$ należy do rzeczywistej wieży kwadratowej. Dlatego $[\Q(u):\Q]$ jest potęgą $2$. Ponieważ $\zeta_n$ spełnia $T^2-uT+1=0$ i nie jest rzeczywisty, mamy $[\Q(\zeta_n):\Q(u)]=2$. Poprzedni lemat daje zatem $\varphi(n)=2[\Q(u):\Q]$, a więc $\varphi(n)$ jest potęgą $2$.

Odwrotnie, jeśli $\varphi(n)$ jest potęgą $2$, to $L=\Q(\zeta_n)$ jest rozszerzeniem Galois z abelową grupą rzędu będącego potęgą $2$. Sprzężenie zespolone generuje podgrupę $H$ rzędu $2$, której ciałem stałym jest $L^H=\Q(u)$: po pierwsze $u$ jest ustalane, po drugie $[L:\Q(u)]=2$. Abeliowość grupy zapewnia, że $\Q(u)/\Q$ jest Galois z abelową grupą $2$-pierwotną. Taka grupa ma szereg podgrup o kolejnych indeksach $2$. Odpowiedniość Galois daje wieżę rozszerzeń kwadratowych od $\Q$ do $\Q(u)$, a wszystkie jej ciała leżą w $\R$. Każdy krok można więc zapisać jako dołączenie dodatniego pierwiastka kwadratowego. Twierdzenie 6.1 pokazuje, że $u$ jest konstruowalne, a zatem również wielokąt.

Pozostaje opisać liczby $n$, dla których $\varphi(n)$ jest potęgą $2$. Jeśli $n=2^r\prod_j p_j^{e_j}$, gdzie $p_j$ są różnymi nieparzystymi liczbami pierwszymi, to
\[
\varphi(n)=\varphi(2^r)\prod_j p_j^{e_j-1}(p_j-1).
\]
Warunek potęgi $2$ wymusza $e_j=1$ i $p_j-1=2^{m_j}$. Jeśli $2^{m_j}+1$ jest pierwsza, to $m_j$ musi być potęgą $2$: gdyby $m_j=2^k d$ z nieparzystym $d>1$, liczba $(2^{2^k})^d+1$ dzieliłaby się przez $2^{2^k}+1$ i nie byłaby pierwsza. Otrzymujemy dokładnie liczby pierwsze Fermata. W drugą stronę każda wypisana w tezie liczba $n$ ma $\varphi(n)$ będące potęgą $2$, co kończy dowód.
\end{proof}

\begin{przyklad}
Liczba $17=2^{2^2}+1$ jest pierwsza, zatem $\varphi(17)=16$ i foremny siedemnastokąt jest konstruowalny. Dowód daje wieżę rozszerzeń kwadratowych; nie wymaga wypisywania pełnego klasycznego wzoru na $\cos(2\pi/17)$.
\end{przyklad}


